% entity-id: axiom-mid
% entity-type: axiom
% lean_manual: true
% system: ZFC
% macro: \TermMid
% args: 0
% notation: \mid

\hypertarget{axiom-mid}{}
\begin{object}[term-mid] % Такое что (Such That)

\textbf{Символ:} $\TermMid{\mid}$

\textbf{Читается как:} «такой, что», «при условии, что».

\textbf{Формальная семантика:}
В теории \Set{множеств} ZFC символ $\TermMid{\mid}$ (или $\colon$) используется в записи \Set{множеств} методом выделения (Set-Builder Notation): $\{x \in A \TermMid{\mid} P(x)\}$ обозначает \Set{множество} всех элементов $x$ из $A$, для которых истинен предикат (свойство) $P(x)$.
Математически выражает отношение отбора или фильтрации элементов.

\textbf{Пример:} $\{x \in \RealNumbers \TermMid{\mid} x > 0\}$
\end{object}